% ============================================================
% Chapter 9: Miscellaneous
% ============================================================
\section{其他知识点}

\subsection{获取系统信息}

\subsection{net\_adm:localhost/0}

\begin{lstlisting}
1> net_adm:localhost().
"centos7-mq1"
\end{lstlisting}

\subsection{监控垃圾回收}

\subsection{清屏}

参考：\href{https://stackoverflow.com/questions/37829916/how-do-i-reset-clear-erlang-terminal}{Stack Overflow - How do I reset/clear erlang terminal}

\subsection{列表打印完全：rp(Term)}

\lstinline|rp(Term)| 使用 shell 已知的 record 定义来打印 term。所有内容都会完整打印，不会像返回值那样限制深度。

\begin{lstlisting}
rp(code:all_loaded()).
\end{lstlisting}

参考：\href{https://www.erlang.org/doc/man/shell.html}{Erlang -- shell}、\href{https://stackoverflow.com/questions/75566211/how-does-it-show-completely-on-erl-terminal-in-erlang}{Stack Overflow}

\subsection{Core Erlang}

\begin{itemize}[leftmargin=1.5em]
    \item \href{https://www.it.uu.se/research/group/hipe/cerl/doc/core_erlang-1.0.3.pdf}{Core Erlang 规范 (PDF)}
    \item \href{https://www.erlang.org/blog/core-erlang-by-example/}{Core Erlang by Example - Erlang/OTP}
    \item \href{https://baha.github.io/intro-core-erlang-2/}{A Gentle Introduction to Core Erlang: Part 2}
    \item \href{https://8thlight.com/insights/the-core-of-erlang}{The Core of Erlang}
\end{itemize}
